\section{Further research questions and topics}
\label{rlc:sec:research}

The following questions seek exact identities or structural
classifications. They are not claims inferred from small numerical
samples. The first four are especially direct continuations of the
proved formulas.

\begin{enumerate}[label=\textbf{Q\arabic*.},leftmargin=2.8em,itemsep=8pt]
\item \textbf{A full multivariate tilt formula.}
Theorem~\ref{rlc:thm:allmoments} evaluates every Taylor coefficient of
$\int e^{\sum_j\lambda_jx_j}\prod_jf_{a,s_j}(x_j)$ under
$\sum x_j=\mu$. Find a compact analytic generating function in all
$\lambda_j$, with its collision limits explicit. The two-factor
series \eqref{rlc:eq:tilted} and the moment-conservation identity supply
nonnegotiable normalization checks.

\item \textbf{Exact unequal-shift mixed order jets.}
Use \eqref{rlc:eq:unequaltransform} to construct a normally convergent
expansion into one central shift, including every independent order
derivative. The periodic analogue already exists in the incoming
unequal-twist report; the real-line version should state its own
admissible contour, convergence domain, and continuation across shift
collisions. Identify which special combinations still have finite
closure without confusing that question with integer partial fractions.

\item \textbf{Odd integral-shift cancellation sets.}
For odd $r$, the common alternating-zeta coefficient in the unit-order
lattice calculation is
$S(\boldsymbol N)=\sum_\nu c_\nu(-1)^{N_\nu}$.
Classify the distinct integer shift sets for which this rational number
is zero. Equal parity is a proved sufficient condition. A classification
would produce additional elementary multivariate integrals in
$\mathbb Q[\pi^2]$. Conversely, $S\ne0$ alone would not prove arithmetic
independence of the surviving odd-zeta combination.

\item \textbf{A central-factorial normal form for moment identities.}
Determine the exact linear relations, as functions of $W,a,\mu$, among
the finite families $\CC_{r+2q,a}(W+d;\mu)$ produced by the moment
algorithm. The product recurrence and covariant ladder already create
relations. A canonical finite reduction would improve symbolic output
without invoking conjectural independence of special values.

\item \textbf{Commensurate dilations of logarithmic arguments.}
Replace profiles by $f_{a_j,s_j}(q_jx)$ for positive rational $q_j$.
Their transforms have cosecant factors with different periods.
Investigate whether a common-denominator decomposition and the Gamma
multiplication formula yield finite Hurwitz families at integer orders,
and determine which all-complex-order sectors retain a finite closure.
The equal-dilation result here is the required base case.

\item \textbf{Complex shift boundaries and explicit contour contacts.}
Continue $a$ beyond the right half-plane while tracking the resolvent
branch point and the poles of the cosecant factor. Derive exact residue
and branch-integral corrections when a chosen contour is crossed.
The ordinary formula for $\Rea a>0$ must remain separate from any new
finite-part convention at the boundary.

\item \textbf{Higher anchored primitives at unequal collisions.}
Construct simultaneous primitives in several independent shifts, with
all constants fixed on a chosen base point and with confluent limits
stated before partial-fraction simplification. The commutator
$[\partial_a,J_a]=-J_a^2$ gives a fully proved one-shift model; the
multishift problem should identify the corresponding mixed correction
operators rather than declare them to commute.

\item \textbf{The same-direction Mellin problem beyond integer orders.}
Return to \eqref{rlc:eq:priorproblem}. Determine a spectral representation
that specializes to the incoming finite double-Hurwitz master but is
holomorphic in both orders. The reciprocal geometry provides a useful
solvable comparison, not a proof that the same finite basis works in the
same-direction geometry.

\item \textbf{The missing Gaussian functional relation.}
Seek a convergent higher-depth or octahedral relation proving the
canonical $S_6$ vector and then the revised $S_8$ vector. The formal
single-product quotient in the manuscript already supplies a separating
functional, so enlarging only that same row span cannot work. A useful
connection to the present Gaussian integral family would have to
produce the missing mixed-color relation explicitly.

\item \textbf{Noninteger convolution depth.}
The contour expression $g(p)^\rho(p+a-1)^{-W}$ can be studied for
noninteger $\rho$ after a branch is fixed. Determine the appropriate
real kernel and exact special-function expression, and identify
precisely why the finite central-factorial product is special to
integer depth. No finite-closure impossibility should be asserted
without specifying the proposed target function class.

\item \textbf{Proof-assistant decomposition and literature reconciliation.}
Formalize the polynomial recurrences, divisibility of $Q_k$, harmonic
coefficient identities, and lattice partial fractions independently of
the analysis. Then isolate weighted Banach-space holomorphy and Fourier
inversion as analytic obligations. Before assigning priority labels,
compare the exact statement set with the three incoming packages whose
contents were not inspected and with the classical convolution
literature.
\end{enumerate}

\section*{Conclusion}
\addcontentsline{toc}{section}{Conclusion}
The reciprocal constraint exposes a total-order multiplier whose
cosecant and resolvent factors yield central-factorial kernels and
Hurwitz order addition. Their interaction gives a finite calculus at
every integer depth and every complex spectral order, including all
logarithmic moments. Stieltjes products, harmonic remainders, Gamma
weights, arithmetic shift collapses, and anchored primitives follow
with their normalization and resonance terms intact.
